 % !TeX root = book.tex
\chapter{What's in it for the Prolog community?}\label{sec:whats-in-it-for-the-prolog-community}

\vspace{0.5cm}

\begin{quotation}
\noindent Coming together is a beginning; keeping together is progress; working
together is success.
%\vspace{0.0cm}
\flushright\textit{Henry Ford}
\vspace{0.5cm}
\end{quotation}

\noindent The earlier chapters developed the Prolog Trinity ecosystem as a technical proposal: Web Prolog as a disciplined profile, Prolog agents as programmable entities, and the Prolog Web as a network of nodes that supports both sequential query interaction and long-lived concurrent conversations. The remaining chapters turn from architecture to value: what, concretely, would this ecosystem make easier or better for the communities that might choose to build with it?

This chapter focuses on the Prolog community itself: implementers, tool builders, teachers, and users who already care about Prolog's future. Chapter~\ref{sec:for-other-communities} widens the lens to neighbouring communities -- logic programming, the Semantic Web, Artificial Intelligence, the Erlang ecosystem, and web programmers -- and asks where the Trinity might offer a useful point of contact. Our aim is to be explicit about trade-offs: which problems the Trinity is meant to address, what it might enable, and which costs it would impose.

%\noindent The journey we have undertaken in this book, exploring the contours of Web Prolog, the potential of Prolog agents, and the vision of a Prolog Web, culminates in a fundamental question: What does the Prolog Trinity ecosystem truly offer? Beyond the technical architecture and the theoretical underpinnings, what are the tangible benefits, the practical advantages, and the opportunities that await individuals and communities who choose to engage with this paradigm? This chapter, as well as the next, are dedicated to answering that precise question. We will peel back the layers of the Prolog Trinity ecosystem and examine its value proposition for a diverse array of stakeholders. This chapter discusses the benefits that the Trinity ecosystem might bring to the dedicated members of the Prolog community, while Chapter~\ref{sec:for-other-communities} takes on the broader spheres of Logic Programming, the Semantic Web, Artificial Intelligence, the Erlang ecosystem, and the vast world of web programmers. Our aim is to articulate not just what the Prolog Trinity \emph{is}, but what it \emph{does} -- what problems it solves, what new horizons it opens, and how it can empower the Prolog community to build a more intelligent, interconnected, and logical web.

For more than five decades, Prolog has served as a testament to the power of logic in computation. Yet, as briefly indicated already in Chapter~\ref{sec:intro}, it has also faced persistent challenges that have, at times, constrained its broader impact. The Prolog Trinity ecosystem is conceived, in no small part, as a direct response to these challenges, offering pathways to reinvigorate the language, and amplify its collective strength.

We will again turn to \textit{Fifty Years of Prolog and Beyond}, where a description of the current state of Prolog is given. We allow ourselves to use the collection of logotypes in Figure~\ref{fig:prolog-systems-logos} as an impressionistic representation of the state of the Prolog ecosystem as it exists today -- the thirteen currently maintained Prolog systems and the communities surrounding them. 

\begin{figure}[h]
    \centering
	\includegraphics[width=9cm]{prolog-systems-logos.png}
    \caption{Logotypes for thirteen Prolog systems: B-Prolog, Ciao, ECLiPSe, GNU Prolog, JIProlog, Scryer Prolog, SICStus Prolog, SWI-Prolog, Tau Prolog, Trealla Prolog, tuProlog, XSB and YAP. Except for Trealla Prolog, these are the systems that FYPB lists as ``currently maintained.'' The matching of systems with logotypes is left as an exercise for the reader.}
    \label{fig:prolog-systems-logos}
\end{figure}

As a background to our discussion, the detailed SWOT analysis presented by the authors of FYPB proves particularly valuable. Their analysis systematically examines the strengths, weaknesses, opportunities, and threats characterizing the Prolog ecosystem, with the aim of reinforcing its strengths, addressing its weaknesses, seizing available opportunities, and mitigating potential threats.

Among the \textit{strengths} that Prolog has, the authors of FYPB not only remind us about the core features of Prolog such as unification and backtracking, but also highlight the many powerful features, such as tabling, concurrency and constraints, that have been developed over the years. They also remind us about how efficient Prolog systems have become. However, for readers familiar with the systems in Figure~\ref{fig:prolog-systems-logos} it should be evident that those strengths are rather unevenly distributed. Some systems offer excellent constraint programming capabilities, others excel at tabling, while others yet focus on concurrency. Some are really fast, others rather slow. XSB's sophisticated tabling mechanisms offer unparalleled performance for certain classes of deductive database queries and logical inferences. ECLiPSe's powerful constraint solvers are indispensable for complex optimization and scheduling problems. SWI-Prolog provides a remarkably comprehensive environment with extensive libraries, making it a popular choice for a wide range of applications. In most comparative speed benchmarks, YAP and SICStus Prolog decisively outperform the other systems.

Occasionally, a new Prolog system emerges that introduces a feature not available in other systems, one that may be difficult or impractical for existing systems to implement. For example, Scryer Prolog and Trealla Prolog have introduced a new form of compact and efficient string representation that emulates a list representation and from the programmer point of view is very much indistinguishable from a list.

As for \textit{threats} and \textit{weaknesses}, the authors of FYPB appear to see \textit{fragmentation} as Prolog's most serious threat. We believe they are right, and the patchwork of logotypes in Figure~\ref{fig:prolog-systems-logos} might be seen as a sign of this threat. Add a user base too small to the problems that Prolog faces -- the language simply is not popular.
%(Another sign is the lack of a system-independent Prolog logotype -- in Chapter~\ref{sec:we-are-the-prolog-web} we will propose one.) 

 %\footnote{Where ``fragmentation'' suggests breaking apart into disconnected or disjointed parts, ``diversification'' implies growth through variety, with an emphasis on expanding into multiple distinct forms, which can suggest a healthy or beneficial development.}


%The fragmentation of the systems landscape as well as the fragmentation of the Prolog community are seen as Prolog's main weakness. 
%
%---

%\noindent 
%We are, however, getting ahead of ourselves. It is prudent first to 

\section{A middle-way approach to systems fragmentation}

According to the authors of FYPB, the fragmentation of the Prolog \textit{systems} landscape is indeed a concern since while most of these systems adhere to the ISO Prolog core standard, many differences appear in the added features that are not handled by ISO. 
%This, of course, leads to code portability problems and difficulties writing text books and tutorials, and might be a threat to further language standardization efforts.

While the diversity is a testament to Prolog's adaptability and the ingenuity within its community, it has historically presented a barrier to collaboration and code reuse. Applications developed in one Prolog system are often difficult to port to another. The traditional pursuit of complete ``portability'' -- the dream of writing Prolog code that runs flawlessly on any system -- while laudable, has proven elusive.
% and, in some ways, perhaps even counterproductive if it were to mean sacrificing the specialized excellences of individual systems for a lowest common denominator.
Fragmentation also leads to difficulties writing textbooks and tutorials, might be a threat to further language standardization efforts, and hinders the development of a healthy ecosystem. 

This section explores three distinct strategies for addressing this long-standing challenge. The first is a bottom-up, pragmatic approach that focuses on incremental improvements and standardization of specific features to gradually enhance portability across systems. In stark contrast, the second approach is a top-down, visionary effort to create a single, unified Prolog that would integrate the best features of all existing systems. Finally, we will examine a third, middle-way approach that champions interoperability over uniformity, proposing a federated model where different Prolog systems can communicate and collaborate in real-time without losing their individual identities. This third approach, which is the focus of this book, is our own proposal.

\subsection{Reducing systems fragmentation one bite at a time}

Since the publication of FYPB, a group of Prolog implementers have begun to address several of the challenges and risks highlighted in the paper. This effort, undertaken over the past few years, involves steps to reduce systems fragmentation. %, and strengthen the community's collective focus.
While it is still an ongoing process, these initial actions appear to represent a growing recognition within the Prolog community of the need to mitigate identified weaknesses and ensure a more cohesive and sustainable future for the language.

\textit{The Prolog Implementers Forum} (with representatives from at least Ciao, ECLiPSe, Logtalk, SWI-Prolog and XSB) is aimed at building consensus on the future of Prolog. In an attempt to address the problem of fragmentation of systems and the code portability problem, forum members have started to try to eliminate some of the differences between systems. Inspired by similar community-driven improvement processes in other programming ecosystems, but tailored to meet the specific needs and challenges of the Prolog community, this has taken the form of \textit{Prolog Improvement Proposals} (PIPs). PIPs are meant to serve as a mechanism for proposing and discussing important enhancements to the Prolog language and its ecosystem. Although not part of the official ISO Prolog standardization effort, they are intended to provide a structured means for documenting, exploring, and evaluating new features, libraries, and extensions that may eventually influence future standardization. 

We do like the idea of using PIPs as a means of improving the portability of as large a fragment of Prolog as possible, thereby helping to reduce fragmentation across systems. This is a \textit{realistic} approach. However, while undoubtedly worthwhile, the PIP effort addresses only one of the questions raised above, namely: How should efforts to increase portability be organized? The chosen answer is: bottom-up, incremental, and in small steps.

PIPs provide a structured arena for implementers to negotiate changes, preventing systems from diverging in isolation. By aligning on specific areas like library interfaces, they shrink the ``portability gap'' and reduce the practical costs of fragmentation. Acting as a bottom-up proving ground for features before potential ISO standardization, PIPs ensure formal standards connect to real-world usage. They do not promise a sweeping vision, but rather a mechanism for practical convergence that keeps diverse systems compatible.

However, this approach faces significant hurdles. Progress is notably slow; despite seven drafted PIPs, no decisions have been reached, and community engagement has stalled. More critically, the effort fails to address the broader question of Prolog's long-term evolution. While the proposals are pragmatic, they lack a unifying vision, and this very absence of direction likely contributes to the sluggish pace. Without a shared roadmap for the future, incremental efforts risk becoming disconnected, piecemeal fixes rather than steps toward a larger goal. Vision acts as a necessary coordinator and motivator, inspiring contribution and helping to prioritize which improvements matter most. When this compelling narrative is missing, discussions fragment, contributors drift away, and even well-designed initiatives lose the momentum required to succeed.

%---
%
%So how can this mitigate the weakness of systems fragmentation? PIPs provide a structured, community-driven process for proposing, documenting, and evaluating changes across systems. This creates a common arena where implementers can collaborate instead of diverging further in isolation. Even if adoption is uneven, the very act of negotiation helps prevent systems from drifting apart entirely.
%
%By ironing out differences in small, well-defined areas (for example, specific library interfaces or language features), PIPs can gradually shrink the ``portability gap'' that makes code reuse and collaboration so hard. This reduces the practical cost of fragmentation for developers and educators.
%
%Although not part of ISO, PIPs can act as a proving ground for features that may later be standardized. In this way, they offer a bottom-up complement to top-down standards processes, ensuring that ISO efforts remain connected to real-world system usage.
%
%In short, while PIPs cannot provide a sweeping vision for Prolog's future, they mitigate fragmentation by creating mechanisms for practical convergence, cultural cohesion, and incremental technical alignment. They do not erase diversity, but they may reduce the risks of Prolog splintering into mutually incompatible islands.
%
%
%There are, however, several problems with this approach. To begin with, progress is slow -- at the time of writing, seven PIPs have been drafted, but no decisions have been made, and the most recent entry in the forum dates back ???? months. More importantly, the effort offers no answer to the broader and more pressing question: How might Prolog and its community evolve in the future? The steps proposed to evolve Prolog as a language are pragmatic, but they can hardly be called visionary. Indeed, although realism is valuable, it may well be that this very lack of vision is itself a contributing factor to the sluggish pace of progress.
%
%A lack of vision can slow progress because it deprives a community of a clear sense of direction and urgency. Without a shared picture of what Prolog could become in five, ten, or twenty years, incremental efforts such as PIPs risk turning into piecemeal technical fixes rather than steps toward a larger goal. Vision serves as both a motivator and a coordinator: it inspires people to contribute, and it provides criteria for deciding which incremental improvements matter most. When vision is absent, discussions easily fragment, priorities become unclear, and contributors may lose interest or drift away. The result is that even well-designed pragmatic initiatives move forward at a sluggish pace, because there is no compelling narrative that binds them together or attracts sustained energy from the community.


%\subsection{One Prolog to Rule Them All}
%\label{sec:one-prolog}
%
%The title of this section is deliberately provocative: the ambition is not that a single Prolog implementation should dominate, but that a single protocol layer should be rich enough to let diverse logic programming systems interoperate as peers on the Web.
%
%Section~\ref{sec:lp-landscape} surveyed the breadth of the logic programming landscape -- from Datalog's guaranteed-terminating queries through WFS-capable systems, probabilistic reasoners, constraint solvers, and multi-paradigm languages such as Picat. Each of these traditions has developed its own execution model, its own tooling, and its own user community, often in relative isolation from one another. The Prolog Web proposes a minimal common interface: any system that can accept a goal (as a Prolog term), compute answers, and return them over HTTP can participate as a node.
%
%What makes this feasible is the profile hierarchy introduced in Chapter~\ref{sec:prolog-agents}. The profiles do not impose a single execution model; they stratify the network by capability. A system that supports only pure, ground-term queries participates at the ISOBASE level. One that adds mutable state and side effects operates at the ISOTOPE level. One that further supports Erlang-style concurrency and process spawning operates at the ACTOR level. This stratification means that no system is required to implement features foreign to its paradigm. A Datalog engine need not pretend to support \texttt{assert} and \texttt{retract}; a probabilistic reasoner need not implement message passing between processes. Each system participates at the level that matches its actual semantics, and clients that issue queries can discover, through the profile, what kind of answers to expect.
%
%The result is not a monoculture but a federation. A network of Prolog Web nodes might include a high-performance Datalog engine for large-scale relational queries, a WFS-capable Prolog system for sound non-monotonic reasoning, a probabilistic logic programming system for uncertainty-aware inference, and a full \textsc{actor}-level node for orchestrating concurrent agents -- all addressable through the same protocol, all discoverable through the same mechanisms, and all able to delegate sub-queries to one another. The phrase ``one Prolog to rule them all'' refers not to a single implementation but to a single architectural pattern -- the Prolog Trinity -- that is flexible enough to serve as a common substrate for the logic programming family as a whole.

\subsection{One Prolog to rule them all}

For an especially visionary perspective, we may turn to one of the contributions in \textit{Prolog: The Next Fifty Years}, where Gupta et al. answer the question -- asked in FYPB but never answered -- of whether it makes sense to aim for a unified language with a resounding ``yes.'' They propose an ambitious vision of a system that seamlessly integrates tabling, constraints, concurrency, and parallelism within a single coherent framework.

%If we want to see just how far a visionary approach can be taken, we need look no further than one of the papers in \textit{Prolog: The Next Fifty Years}. There, Gupta and his co-authors answer the question “Does it make sense to aim for a unified language?” with an emphatic yes, and sketch a bold design that brings together tabling, constraints, concurrency, and parallelism in a single system.

%\noindent For an extremely visionary approach we turn to one of the papers in PN50Y, where Gupta et al. seems to answer the question ``Does it make sense to aim for a unified language?'' with a resounding ``Yes!'' and propose a particularly ambitious vision of a language that seamlessly integrates tabling, constraints, concurrency, and parallelism within a single system. 

The idea of a unified system is not new. In a 2016 discussion on Stack Overflow,\footnote{\url{https://stackoverflow.com/questions/35668475/prolog-vs-erlang-and-other-functional-languages}} Jan Wielemaker proposed:

\begin{quote}
Ideally we should sit together and assemble a new system that combines the speed of YAP with the tabling of XSB, the constraints of ECLiPSe, and the environment and interfaces of SWI-Prolog.
\end{quote}

\noindent Let us give this hypothetical system a name. Interestingly, by permuting the initial letters of the component systems' names, one arrives at a playful and memorable moniker for the new platform: SEXY Prolog!

The approach of Gupta et al goes even further than Wielemaker though, as they also aim to incorporate \textit{Answer Set Programming} (ASP) into Prolog, or more specifically s(CASP), which combines the strengths of ASP and constraint logic programming in a top-down, query-driven model of computation that includes support for explanation generation. The hope is that this leads to a sophisticated evolution that consolidates various extensions and paradigms into a cohesive programming environment even capable of commonsense reasoning. This integration is claimed to be designed not only to expand Prolog's expressive power but also to align it more closely with the evolving demands of artificial intelligence. % and other advanced applications. %For the purpose of this book, and to emphasize the importance of standardization, we shall refer to this unified language as ISO Prolog 2.0.


So what happened to SEXY Prolog? Its spirit of integration lives on, and SWI-Prolog has been repeatedly described as a system that ``glues together'' many Prolog research and application innovations into a broadly usable environment. In recent years, SWI-Prolog has incorporated advanced features such as CLP (inspired by ECLiPSe -- the prototypical CLP-oriented Prolog system), tabling (a key feature of XSB -- but in an implementation that according to the manual is a ``first prototype'' requiring further enhancement to rival a mature system such as XSB.''), and for speed, Just-In-Time indexing (for which the manual says  that ``YAP's indexing has been the inspiration for enhancing SWI-Prolog's indexing capabilities.''). On top of this, SWI-Prolog even has an implementation of s(CASP).\footnote{https://github.com/SWI-Prolog/sCASP}

In this way, users do not need to switch between multiple systems for different features; they can stay in SWI-Prolog and get ``good enough'' versions of everything.
%While there is no publicly documented `ECLiPSe -> SWI-Prolog' borrow, the historical role of ECLiPSe as a prototypical solver-oriented Prolog system, combined with SWI-Prolog's integrative design philosophy and its adoption of standard CLP machinery (attributed variables, multiple CLP solver libraries), make it plausible that ECLiPSe served at least as a reference model or informal inspiration for SWI-Prolog's CLP infrastructure.
%
%Furthermore, SWI-Prolog has incorporated tabling, a key feature of XSB. However, the SWI-Prolog manual modestly describes its tabling implementation as a ``first prototype'' requiring further enhancement to rival a mature system such as XSB. %The key takeaway is not the creation of a monolithic new Prolog, but the organic evolution of existing systems to embrace powerful features from across the Prolog landscape.
%
%As far as speed goes, some of YAP's optimizations influenced SWI-Prolog. For example, YAP provides full JIT indexing, including indexing arguments of compound terms, and according to the SWI-Prolog manual, ``YAP's indexing has been the inspiration for enhancing SWI-Prolog's indexing capabilities.''
However, it seems that for Gupta et al, ``good enough'' is, in fact, not good enough. 
%They emphasize that the added advanced features must be \textit{seamlessly} integrated, 
What matters is not merely the presence of advanced features, but their \textit{seamless} integration -- a standard that, in their view, SWI-Prolog probably does not yet meet. They therefore contend that research ``to realize a system where all these advanced features are efficiently and seamlessly integrated and available in a single system must continue.'' 

%PN50Y: Unfortunately, various strategies developed to make execution more efficient are not all available in a single Prolog system where they work seamlessly with each other. We believe that future research in Prolog must focus on building a unified system that supports tabling, constraints, coroutining, parallelism, and concurrency [27]. These enhancements must also work seamlessly with each other.
%
%PN50Y: Research to realize a system where all these advanced features are efficiently and seamlessly integrated and available in a single system must continue. Our hope is that such a logic programming system will be realized in the future.

The approach of Gupta et al is extremely interesting, but we suspect that it is not likely to result in anything practically useful any time soon. It would likely take decades for a unified language of this kind to materialize and become useful. The question is also whether the result, if it materializes, will still be Prolog? We believe it is more likely to end up as another Mercury, Oz or Picat. 
%, rather than as ISO Prolog 2.0. 
They are excellent programming languages, but they are not Prolog. % Nevertheless, unification provides a visionary pole, a ``north star'' that defines an ideal (even if unreachable).

%And it does not reduce fragmentation... or does it? Instead of each system choosing its own subset of features, a unifying semantics allows one system to cover all use cases soundly and elegantly. This would reduce theoretical fragmentation: fewer incompatible dialects or semantics; more shared research direction. Offers a long-term research roadmap to avoid Prolog drifting into isolated niches.
%
%Together, these two approaches cover both the engineering pathway and the theoretical foundation. Without Wielemaker, Prolog risks theoretical beauty but practical obscurity. Without Gupta et al., Prolog risks practical adoption but a messy, inconsistent semantics. Together, they point to a less fragmented, more coherent Prolog ecosystem.

\subsection{The middle way: interoperability as a federating force}

The history of programming language standardization is largely a
history of unification attempts: efforts to define a single language,
a single type system, a single runtime, or a single ontology that
all participants must adopt. Such efforts have produced valuable
results -- ISO Prolog, the W3C Semantic Web stack, and the JVM are
notable examples -- but they also carry a characteristic risk.
Unification demands consensus on internal design decisions; the
broader and deeper the consensus required, the harder it is to
achieve and maintain. In the Prolog world specifically, the
fragmentation of implementations and the difficulty of converging on
shared standards have been persistent challenges.

The Trinity takes a different approach. It defines shared
\emph{interaction surfaces} -- profiles, protocols, and APIs -- while
leaving internal implementation choices free. A node need not run a
particular Prolog system; it need only conform to a profile that
specifies which predicates, transports, and behavioural contracts it
supports. Different nodes may implement the same profile in different
ways, with different optimisations, different libraries, and
different internal architectures. What is standardised is the
\emph{interface}, not the \emph{engine}.

This is federation in a precise sense: independent parties agree on
how to communicate without agreeing on how to think. The agreement is
operational (we will exchange messages in this format, over these
transports, with these semantics) rather than ontological (we will
represent the world in the same way). This is also why the Trinity
can position itself as scaffolding around existing Prolog systems
rather than as a replacement for any of them.

Figure~\ref{fig:prolog-web-intro-200} illustrates the concept in
concrete terms. Rather than collapsing the distinctive strengths of
SWI-Prolog, ECLiPSe, XSB, and YAP into a single uniform system, we
envision an ecosystem where each system is wrapped in a node
implementation that allows it to converse with others. The diversity
of approaches is preserved, yet they can all contribute to a common
fabric of cooperation. While participation is voluntary, a node
interface of this kind is indispensable for any system aspiring to
take part in this shared Prolog Web.

\begin{figure}[h]
\centering
\includegraphics[width=11cm]{prolog-web-intro-200.png}
\caption{A federated middle way: interoperability without forcing a
single unified system.}
\label{fig:prolog-web-intro-200}
\end{figure}

As demonstrated in the proof-of-concept implementations
accompanying this book, Web Prolog nodes owned by different entities
and running in different locations can offload computations to one
another and engage in mutual communication. This represents a
significant step towards real-time interoperability. However, the
discussion so far has been limited to homogeneous cases where nodes
run on the same underlying Prolog system. The next challenge is to
extend this capability to heterogeneous settings, where nodes may
rely on different Prolog systems yet must still achieve the same
level of seamless, real-time interoperability.


\subsubsection{Portability vs.\ real-time interoperability}

Much like the work pursued by the Prolog Implementers Forum (PIF),
the advancement of this middle way requires cooperation among system
implementers. However, while the PIF emphasises \emph{portability}
-- the ability to move code between platforms offline -- our proposal
extends this focus to \emph{real-time interoperability}. Portability
enables code to move between systems; interoperability enables the
systems themselves to interact during execution.

Because runtime conformance issues cannot be easily fixed on the fly,
the demands placed on Web Prolog systems must be strictly defined.
The Prolog Trinity ecosystem offers a vision of \emph{cross-system
interoperability as a federating force}. Instead of striving to make
all Prolog systems behave identically, Web Prolog functions as a
\textit{lingua franca} that enables distinct systems to communicate,
collaborate, and interoperate seamlessly.


\subsubsection{Reframing fragmentation as diversification}

We emphasise that this approach allows Prolog systems to
\textit{retain their differences} on many levels. We encourage a
``let a thousand flowers bloom'' philosophy: implementers are free to
add extensions and optimisations, provided they do not interfere with
the Web Prolog layer. Through calls to the host language, a Prolog
Web node can still offer the functionality of these unique
extensions; indeed, their presence might be exactly why a developer
chooses that specific node implementation.

Imagine a future where an application orchestrated by Web Prolog can
dispatch a deductive query to a remote XSB engine, feed the results
into an ECLiPSe solver, and manage the user interface via
SWI-Prolog. Each system contributes its peak capability, forming a
federated ecosystem where developers do not have to be content with
``good enough'' versions of everything, but can aim for ``the best
there is.''

From this vantage point, the current landscape of the Prolog
community changes. What is often labelled as fragmentation can be
reinterpreted in a more constructive light: as \emph{diversification}.
This shift in framing underscores the presence of a dynamic, albeit
loosely coupled, constellation of systems that collectively sustain
innovation and resilience.


\subsubsection{Engineering a collaborative future}

Much like the grand vision for a Unified Prolog, our approach entails
substantial work. However, it demands less in terms of foundational
research and more in terms of disciplined \textit{engineering}. We
already possess mature building blocks that need to be consolidated
rather than re-invented. Furthermore, \textit{inclusivity} is a core
principle. To ensure that the barrier to entry is not insurmountably
high -- particularly for systems lacking multi-threading -- the
profile hierarchy proposed in Chapter~\ref{sec:prolog-agents} allows
even the least powerful Prolog systems to integrate by conforming to
simpler profiles.

Ultimately, the viability of this middle way depends on a symbiotic
relationship with the Prolog Implementers Forum. The PIF's efforts to
catalogue divergences and clarify semantics form the substrate upon
which our interoperability architecture is built. In return, the
Trinity ecosystem acts as a unifying testbed: a practical arena where
abstract agreements take on technical reality. This offers a gradual
path toward unification, where interoperability precedes
standardisation, allowing the Forum's work to crystallise into a
shared, operationally validated Prolog core that systems can adopt
without sacrificing their individuality.

If the Prolog Web succeeds, it will not be because every Prolog
system has been replaced by a single implementation, but because many
implementations have adopted compatible interaction surfaces. Success
is measured by interoperation and reuse, not by replacement. The same
principle extends to neighbouring communities: the Trinity does not
ask the Erlang community to adopt Prolog, or the Semantic Web
community to abandon RDF. It asks whether there is value in a shared
execution substrate for logic-based agents on the Web -- and offers a
concrete proposal for what that substrate might look like.



%\subsection{The middle way: interoperability as a federating force}
%
%In our approach to mitigating the threat of systems fragmentation, and drawing on the capabilities of the Prolog Trinity ecosystem, we propose a \textit{middle way}: a path situated between the cautious, realistic approach of the Prolog Implementers Forum and the grand vision of a Unified Prolog ``to rule them all.''
%
%Figure~\ref{fig:prolog-web-intro-200} illustrates the concept. Rather than collapsing the distinctive strengths of SWI-Prolog, ECLiPSe, XSB, and YAP into a single uniform system, we envision an ecosystem where each system is wrapped in a node implementation that allows it to converse with others. In this way, the diversity of approaches is preserved, yet they can all contribute to a common fabric of cooperation. While participation is voluntary, a node interface of this kind is indispensable for any system aspiring to take part in this shared Prolog Web.
%
%\begin{figure}[h] 
%\centering 
%\includegraphics[width=11cm]{prolog-web-intro-200.png} 
%\caption{A federated middle way: interoperability without forcing a single unified system.} 
%\label{fig:prolog-web-intro-200} \end{figure}
%
%\noindent As demonstrated in the proof-of-concept implementations accompanying this book, Web Prolog nodes owned by different entities and running in different locations can offload computations to one another and engage in mutual communication. This represents a significant step towards real-time interoperability. However, the discussion so far has been limited to homogenous cases where nodes run on the same underlying Prolog system. The next challenge is to extend this capability to heterogeneous settings, where nodes may rely on different Prolog systems yet must still achieve the same level of seamless, real-time interoperability.
%
%\subsubsection{Portability vs. real-time interoperability}
%
%Much like the work pursued by the Prolog Implementers Forum (PIF), the advancement of this middle way requires cooperation among system implementers. However, while the PIF emphasizes \emph{portability} -- the ability to move code between platforms offline -- our proposal extends this focus to \emph{real-time interoperability}. Portability enables code to move between systems; interoperability enables the systems themselves to interact during execution.
%
%Because runtime conformance issues cannot be easily fixed on the fly, the demands placed on Web Prolog systems must be strictly defined. The Prolog Trinity ecosystem offers a vision of \emph{cross-system interoperability as a federating force}. Instead of striving to make all Prolog systems behave identically, Web Prolog functions as a \textit{lingua franca} that enables distinct systems to communicate, collaborate, and interoperate seamlessly.
%
%\subsubsection{Reframing fragmentation as diversification}
%
%We emphasize that this approach allows Prolog systems to \textit{retain their differences} on many levels. We encourage a ``let a thousand flowers bloom'' philosophy: implementers are free to add extensions and optimizations, provided they do not interfere with the Web Prolog layer. Through calls to the host language, a Prolog Web node can still offer the functionality of these unique extensions; indeed, their presence might be exactly why a developer chooses that specific node implementation.
%
%Imagine a future where an application orchestrated by Web Prolog can dispatch a deductive query to a remote XSB engine, feed the results into an ECLiPSe solver, and manage the user interface via SWI-Prolog. Each system contributes its peak capability, forming a ``federated ecosystem'' where developers do not have to be content with ``good enough'' versions of everything, but can aim for ``the best there is.'' 
%
%From this vantage point, the current landscape of the Prolog community changes. What is often labeled as fragmentation can be reinterpreted in a more constructive light: as \emph{diversification}. This shift in framing underscores the presence of a dynamic, albeit loosely coupled, constellation of systems that collectively sustain innovation and resilience.
%
%\subsubsection{Engineering a collaborative future}
%
%Much like the grand vision for a Unified Prolog, our approach entails substantial work. However, it demands less in terms of foundational research and more in terms of disciplined \textit{engineering}. We already possess mature building blocks that need to be consolidated rather than re-invented. Furthermore, \textit{inclusivity} is a core principle. To ensure that the barrier to entry is not insurmountably high -- particularly for systems lacking multi-threading -- the profile hierarchy proposed in Chapter~\ref{sec:prolog-agents} allows even the least powerful Prolog systems to integrate by conforming to simpler profiles.
%
%Ultimately, the viability of this middle way depends on a symbiotic relationship with the Prolog Implementers Forum. The PIF's efforts to catalogue divergences and clarify semantics form the substrate upon which our interoperability architecture is built. In return, the Trinity ecosystem acts as a unifying testbed: a practical arena where abstract agreements take on technical reality. This offers a gradual path toward unification, where interoperability precedes standardization, allowing the Forum's work to crystallise into a shared, operationally validated Prolog core that systems can adopt without sacrificing their individuality.


%\subsection{The middle way: interoperability as a federating force}
%
%In our approach to mitigating the threat of systems fragmentation, and drawing on the capabilities of the Prolog Trinity ecosystem, we propose a middle way: one that lies \textit{between} the cautious and realistic path taken by the Prolog Implementers Forum and the grand vision of a Unified Prolog ``to rule them all.''
%
%Figure~\ref{fig:prolog-web-intro-200} illustrates the idea. Rather than collapsing the distinctive strengths of SWI-Prolog, ECLiPSe, XSB, and YAP into a single uniform system, we envision an ecosystem in which each system is wrapped in a node implementation that allows it to converse with the others. In this way, the diversity of approaches is preserved, yet they can all contribute to a common fabric of cooperation. Participation is of course voluntary, but for a system to take part in this shared Prolog Web, a node interface of this kind is indispensable.
%
%\begin{figure}[h]
%    \centering
%	\includegraphics[width=11cm]{prolog-web-intro-200.png}
%    \caption{The sexy Prolog Trinity: End-user clients in interaction with sexy Prolog agents on the sexy Prolog Web. [TODO: Perhaps VR headset and Furhat robot should be remeved?]}
%    \label{fig:prolog-web-intro-200}
%\end{figure}
%
%As explained in earlier chapters, and as demonstrated by the proof-of-concept implementation accompanying this book, Web Prolog nodes with different owners and running at different locations can offload computations to each other and engage in mutual communication. This represents an important step towards real-time interoperability. However, the discussion so far has been limited to cases where both nodes run on the same underlying Prolog system. 
%
%The next challenge is to extend this capability to heterogeneous settings, where nodes may rely on different Prolog systems -- or even different programming languages altogether -- yet must still achieve the same level of seamless, real-time interoperability.
%
%
%\subsubsection{Portability, interoperability and Web Prolog as a lingua franca}
%
%Like the work pursued by the Prolog Implementers Forum, the advancement of this middle way can only be achieved through cooperation among system implementers. Its success depends on the bottom-up, practical collaboration between those responsible for maintaining and evolving Prolog systems. Work on the Trinity Prolog ecosystem must be carried out top-down.
%
%Note that the emphasis placed by the Prolog Implementers Forum on \emph{portability} aligns closely with our own priorities. The ability to share programs, tools, and components across different Prolog implementations is as essential to our vision of a way to reduce systems fragmentation as it is to theirs. 
%
%---
%
%Porting programs between platforms takes place offline, and with a bit of work problems can usually be taken care of. Real-time interoperability is a runtime phenomenon, and conformance issues arising at runtime cannot be easily fixed. Thus, the conformance demands that we must place on Web Prolog systems need to be very strict.
%
%---
%
%However, our proposal extends this focus by treating portability as a necessary foundation for \emph{instant portability} and \emph{real-time interoperability} \textendash{} that is, the ability for distinct systems to communicate and cooperate effectively across distributed environments. In this sense, portability enables code to move between systems, while interoperability enables systems themselves to interact in the construction of a coherent, distributed infrastructure for logic programming on the Web.
%
%The Prolog Trinity ecosystem, and Web Prolog in particular, offers a pragmatic yet powerful vision: \emph{cross-system interoperability as a federating force}. Instead of striving to make all Prolog systems behave identically, Web Prolog is designed to function as a \textit{lingua franca} that enables these distinct systems to communicate, collaborate, and interoperate seamlessly across the Web. 
%
%
%\subsubsection{Let a thousand flowers bloom!}
%
%We note (again) that the approach we suggest will allow Prolog systems to \textit{stay incompatible} on many levels and in many ways, as long as they remain compatible on the level of Web Prolog. All sorts of extensions to be made to a Prolog system are allowed (and appreciated -- let a thousand flowers bloom!), as long as they do not interfere with the Web Prolog layer. Through calls to the host language, a Prolog Web node based on such a platform will still be able to offer the functionality of such extensions and their presence might be what makes a developer choose it over other node implementations. 
%
%Imagine a future where your application, written in Web Prolog, can be instantly ported to a node implemented on top of another system, or imagine that if orchestrated by Web Prolog, it can dispatch a complex deductive query to a remote XSB engine, take the resulting knowledge and feed it into an ECLiPSe solver to optimize resource allocation under intricate constraints, and then utilize an SWI-Prolog instance to manage the web interface and user interaction. Each system contributes its peak capability, forming a ``federated ecosystem'' where the whole is far greater than the sum of its parts.
%
%No longer are developers confined to the capabilities of a single chosen Prolog implementation. Instead, they can architect distributed applications that strategically leverage the best-in-class features from across the entire Prolog landscape, fostering innovation and enabling the development of sophisticated and powerful intelligent systems.
%
%
%\subsubsection{Less research, more engineering}
%
%Much like the grand vision for a Unified Prolog, our approach entails substantial work in terms of design and implementation. However, in contrast to that more ambitious objective, our proposal demands much less in terms of foundational research. We already offer mature building blocks that need to be consolidated and integrated rather than fundamentally re-invented. What is primarily required is disciplined engineering rather than the pursuit of new theoretical breakthroughs -- engineering efforts that extend established techniques rather than necessitating research into the foundations of logic programming.
%
%
%\subsubsection{Inclusivity}
%
%It is important to ensure that the barrier to entry is not insurmountably high, so \textit{inclusivity} is a core principle of the Prolog Trinity ecosystem. Ideally, all thirteen systems presented at the start of this chapter will be able to participate. However, for systems that lack multi-threading, implementation becomes much harder, so realistically, not all Prolog systems will implement full actor-based concurrency. To ensure as much inclusivity as possible, the profile hierarchy proposed in Chapter~\ref{sec:prolog-agents} ensures that even the least powerful among Prolog systems can be integrated as nodes by conforming to the appropriate profile within this hierarchy. This model not only maximizes participation across the Prolog community but also promotes a flexible system where varying degrees of functionality coexist harmoniously.
%
%
%\subsubsection{A change of perspective}
%
%In discussions of programming language ecosystems, fragmentation is frequently regarded as a liability. It evokes images of disorganization, redundant effort, and a diminished capacity to attract users or exert influence. Yet, in the case of the contemporary Prolog ecosystem, what is often labeled as fragmentation may instead be interpreted in a more constructive light -- as evidence of organic evolution, pluralism in design approaches, and adaptive responses to varied use cases and implementation contexts. Viewed from this angle, the current landscape might be more appropriately characterized not as one of fragmentation, but of \emph{diversification}. This shift in framing underscores the presence of a dynamic, albeit loosely coupled, constellation of systems and practices that collectively sustain innovation and resilience within the broader Prolog community.
%
%\subsubsection{Where the middle way makes a difference}
%
%The viability of any middle-way approach rests on more than technical ingenuity; it depends on a shared commitment among implementers to converge, incrementally, on a set of common foundations. In this respect, the role of the Prolog Implementers Forum (PIF) is indispensable. Its ongoing efforts to catalogue divergences, clarify intended semantics, and maintain a living record of consensus points form the very substrate upon which a federated, interoperability-first architecture must be built. Without such groundwork -- without a stable, community-curated understanding of what \textit{Prolog} currently is -- there is little hope of constructing a robust fabric allowing heterogeneous systems to speak to one another in predictable ways.
%
%At the same time, the Prolog Trinity ecosystem can offer something back. By proposing a concrete architecture for real-time interoperability -- complete with node profiles, communication protocols, and executable specifications -- it provides a practical arena in which the PIF's abstract agreements can take on technical reality. In this sense, the Trinity acts as a unifying testbed: a place where implementers can verify compatibility, observe the friction points between systems, and gain empirical guidance about which forms of unification are within reach. The vision of a future unified Prolog has often seemed either too ambitious or too premature; what the middle-way approach offers is a gradual path toward that vision, one where interoperability precedes unification, but also gently prepares the ground for it. If successful, it may allow the Forum's work to crystallise not merely into a set of recommendations but into a shared, operationally validated Prolog core -- one that multiple systems can adopt without sacrificing their individuality.


\section{How to deal with community fragmentation}

\begin{quotation}
\noindent The Web is more a social creation than a technical one. I designed it for a social effect -- to help people work together -- and not as a technical toy. The ultimate goal of the Web is to support and improve our weblike existence in the world.
\vspace{-0.3cm}
\flushright \textit{Tim Berners-Lee, 2000}
\vspace{0.1cm}
\end{quotation}

\noindent In addition to systems fragmentation, the authors of \textit{Fifty Years of Prolog} (FYPB) identify the fragmentation of the Prolog \textit{community} as a significant concern for both the standardization process and the long-term growth of the language. This social fragmentation mirrors the technical landscape: application developers and researchers congregate around specific systems (SWI-Prolog, Ciao, XSB, etc.), forming distinct sub-communities.

While these groups foster valuable local discussions, their isolation leads to duplicated efforts, a dilution of Prolog's collective voice in the wider technological discourse, and a slower dissemination of innovations. This disjointed state hampers collaboration and complicates efforts to maintain a unified language standard.

We argue that by utilizing the Prolog Trinity ecosystem, we can bridge these sub-communities through four key pillars: interaction, collaboration, learning, and shared identity.

\subsection{Coming together through interaction}

A programming language ecosystem, defined as a dynamic network of software, developers, and system implementers, relies on continuous interaction for its health and sustainability. The advent of the World Wide Web fundamentally transformed these ecosystems, shifting them from localized groups to interconnected global communities via platforms like GitHub and Stack Overflow. While the Prolog community has leveraged these tools to some extent, the authors of \textit{Fifty Years of Prolog} highlight the absence of a strong, unified online presence compared to the vibrant communities surrounding languages like Python or JavaScript.

The ``All Things Prolog'' website, sponsored by the Association for Logic Programming, represents the first significant step toward a system-independent homepage.\footnote{\url{https://prolog-lang.org/}} However, to truly combat fragmentation, Prolog requires more than a static homepage; it needs a dynamic \textit{portal}. We must distinguish between the \textit{technical} Prolog Web -- the infrastructure of nodes and agents -- and the \textit{social} Prolog Web -- the network of users and developers. These two concepts should not exist in isolation. In the spirit of ``We Are the Web,'' we envision a portal that integrates them, serving as the definitive gateway to the Prolog Trinity ecosystem.

Such a portal would be far more than a repository of links. It would serve as a living library, hosting comprehensive documentation and tutorials specifically designed to guide users through Web Prolog and agent building. Beyond static resources, it would function as a dynamic registry, providing a service directory of available public Prolog nodes and services, making the distributed nature of the ecosystem visible and accessible. To foster human connection, the portal would act as a community hub featuring forums, news, and collaborative tools, ensuring that the social layer is as robust as the technical one. Finally, it would host the playground (see Section~\ref{sec:playgrounds}), an interactive environment for instant experimentation, allowing visitors to write and run code without leaving the browser.

The underlying philosophy is that technical interoperability can be a powerful catalyst for community cohesion. If we can make the many different Prolog systems talk to each other over the Web, there is hope that the Prolog sub-communities will start talking to each other too.

\subsection{Coming together through collaboration and sharing}

Collaboration is only made possible through interaction. To strengthen the community, we propose a culture of ``dogfooding'' -- where the community actively uses the Prolog Web to build the Prolog Web. The most immediate early adopters should be Prolog users and implementers themselves.

This extends to the concept of \textit{sharing}. Currently, we share code and ideas. The Prolog Trinity ecosystem enables us to share \textit{computational resources} and \textit{live processes}. Imagine a future where a developer can not only share a library on GitHub but also spawn a live demonstration agent on a public node, allowing others to interact with it immediately. This shifts the paradigm from social coding to \textit{social computing}.

FYPB notes that the community lacks a standard package manager and that portability issues extend to libraries and tools such as testing and documentation infrastructure. This strengthens the case for unifying web-based environments at the node level: a shared, profile-governed Web Prolog node can provide a single distribution channel for code, a single execution target for teaching material, and a coherent place to attach tool support, while still allowing multiple underlying implementations.


\subsection{Coming together through learning}

\begin{quotation}
\noindent Ah, Prolog. Sometimes spectacularly smart, other times just as frustrating. You'll get astounding answers only if you know how to ask the question. Think \textit{Rain Man}. I remember watching Raymond, the lead character, rattle off Sally Dibbs' phone number after reading a phone book the night before, without thinking about whether he should. With both Raymond and Prolog, I often find myself asking, in equal parts, ``How did he know that?" and ``How didn't he know that?" He's a fountain of knowledge, if you can only frame your questions in the right way.
\vspace{-0.3cm}
\flushright\textit{Bruce Tate}
\vspace{0.5cm}
\end{quotation}


\noindent The future of any technology is inextricably linked to education. For Prolog to thrive and expand its reach, we must provide engaging, accessible, and modern learning experiences. The inherently interactive and distributed nature of the Prolog Trinity, particularly Web Prolog, offers an unprecedented opportunity to create ``Prolog Playgrounds'' that go far beyond traditional textbook exercises or standalone interpreters.


\subsubsection{Prolog playgrounds}\label{sec:playgrounds}

Figure~\ref{fig:prolog-playgrounds} shows three of the currently existing web-based Prolog playgrounds: SWISH (based on SWI-Prolog), the Ciao Prolog playground, and the Tau Prolog sandbox. These tools allow visitors to program in Prolog directly in the browser without installing software, and they have proven invaluable for education and quick prototyping.

\begin{figure}[h]
    \centering
    \includegraphics[width=\textwidth]{prolog-playgrounds.png}
    \caption{Current web-based environments: SWISH (SWI-Prolog), Ciao Prolog playground, and the Tau Prolog sandbox.}
    \label{fig:prolog-playgrounds}
\end{figure}

However, as things now stand, these playgrounds unfortunately mirror the very fragmentation of the community they serve. While they share the goal of bringing Prolog to the Web, they do not support the \textit{same} language. Each playground runs a specific dialect -- SWI, Ciao, or a JavaScript-based ISO subset -- meaning that code written in one environment may not be portable to another without modification. A student or developer learning via SWISH encounters friction when moving to the Ciao playground, not only due to syntax divergences but also because the GUIs behave differently. The shortcuts, query mechanisms, and visualization tools are inconsistent, forcing users to relearn the tool rather than focusing on the logic. Furthermore, these distinct implementations represent a significant amount of duplicated engineering effort. Significant resources have been poured into developing three separate, siloed web interfaces -- effort that could have been far more effective if directed towards a collaborative initiative.

Fortunately, it is not too late to change course. Building a unified Prolog playground served directly by a Web Prolog node is entirely feasible. An ISOTOPE node would be fully capable of serving a standard Prolog environment suitable for general instruction. Moreover, if served by an ACTOR node, the environment could evolve into a playground for teaching Erlang-style actor programming in Web Prolog, empowering a new generation of developers to experiment with distributed concurrency within a single, coherent framework.

Prolog is indeed known to have a steep learning curve, and Erlang is probably not \textit{that} far behind. Alas, there is nothing in Web Prolog that promises to make its learning curve less steep, but at least, to be sure, people already familiar with both Prolog and Erlang will have an easy task picking up Web Prolog, including the trickier areas such as concurrent programming. Prolog programmers without Erlang experience will be able to use Web Prolog out of the box -- they only need to learn new things if they want to delve into concurrent programming in Web Prolog, and there are Erlang textbooks from which the principles as well as some of the practice can be learned. Erlang programmers not familiar with Prolog will have some new things to learn, but the close affinities between the two languages are likely to be of help here too.

By making the learning process more interactive, more contextualized within modern web technologies, and more directly connected to the exciting possibilities of distributed intelligent systems, these Prolog Playgrounds can ignite the curiosity of a new generation of programmers and researchers. They can showcase Prolog not as a historical artifact, but as a living, breathing language uniquely suited to addressing some of the most exciting challenges in computing today.


\subsubsection{Learning materials}\label{sec:books}

Staying close to traditional Prolog also makes it possible to reuse some of the old textbooks written for teaching Prolog. There are plenty of really good Prolog textbooks around, and many can be downloaded for free. Some great books teaching Erlang exist too -- and again, some of them are free.

\begin{figure}[h]
    \centering
	\includegraphics[width=11cm]{books}
    \caption{Textbooks for learning Prolog and Erlang -- there are many others!}
    \label{fig:books}
\end{figure}

\noindent Learning concurrent and distributed programming in Web Prolog by reading Erlang books is surprisingly easy. Of course, at some point is would be wise to produce a dedicated book or at least a tutorial.


\subsection{Coming together through rallying around a shared symbol}

\begin{quotation}
\noindent One of the deep mysteries to me is our logo, the symbol of lust and knowledge, bitten into, all crossed with the colors of the rainbow in the wrong order. You couldn't dream a more appropriate logo: lust, knowledge, hope and anarchy
\vspace{-0.3cm}
\flushright  \textit{Jean-Louis Gass\'ee}
\vspace{0.3cm}
\end{quotation}

\noindent A logotype gives a programming language a unique look and feel that makes it instantly recognizable and easier to remember. It is more than just a nice design -- it is a way for users to connect, feel part of a community, and rally around a shared symbol. The logo also says something about what the language stands for, helping people get a quick sense of its vibe. Plus, it is useful for spreading the word; a good logo can make the language more visible in tutorials, online discussions, and conferences, making it stick in people's minds and building a legacy over time. Here are the logotypes for four well-established languages:

\begin{figure}[h]
    \centering
	\includegraphics[width=9cm]{pl-logos}
    \caption{Logotypes for four programming languages.}
    \label{fig:pl-logos}
\end{figure}

\noindent By contrast, there is no Prolog logo broadly recognized across sub-communities. It may be that lack of a formal logo is intended to make a statement, signaling that Prolog is academic software? However, in recent years, even niche languages have started adopting logos as branding and community-building become more central, especially in open-source and web-facing contexts. We do believe it is time for Prolog to follow suit. 

Fortunately, we need not begin from scratch. The field of logic programming already possesses an emblematic resource: the white-on-black geometric symbol at the heart of the logo for the \textit{Association for Logic Programming} (ALP).

\begin{figure}[h]
    \centering
	\includegraphics[width=8cm]{alp-logo}
    \caption{Logotype for the Association for Logic Programming.}
    \label{fig:alp-logo}
\end{figure}

\noindent %At the heart of the ALP logo lies a small, subtle white-on-black emblem. 
Familiar to logicians, the three glyphs $\vDash$, $\equiv$ and $\vdash$ stand for semantic entailment, logical equivalence and syntactic inference, respectively. Though underutilized and perhaps underrecognized, the equation expresses a core thesis of logic programming: that syntax and semantics, process and meaning, are formally intertwined. The equation, while not fully realized in practical Prolog, remains a guiding principle of logic programming -- and a fitting philosophical emblem for Prolog.

In the spirit of reuse and renewal, and adding a frame and a bit of color, we present the following proposal for a Prolog logotype:

\begin{figure}[h]
    \centering
	\includegraphics[width=4cm]{prolog-logo}
    \caption{A possible Prolog logotype.}
    \label{fig:prolog-logo}
\end{figure}

\noindent By replacing the name ``Prolog'' with the names of other logic programming languages, and varying the color of the name and the frame, logotypes for other languages implementing this paradigm can be formed. Figure~\ref{fig:lp-languages-logos} shows four examples:

\begin{figure}[h]
    \centering
	\includegraphics[width=11cm]{lp-languages-logos}
    \caption{Logotypes for four other logic programming languages.}
    \label{fig:lp-languages-logos}
\end{figure}

\noindent Of course, the logos need to look 2025-modern, not 1995-modern, so these visual mockups need to be professionally rendered if they are to be taken seriously.

Perception matters. A modern, compelling visual identity -- a Prolog Web logotype -- can serve as a powerful symbol of unity, innovation, and shared purpose. It helps to present Prolog not as a relic of the past, but as a forward-looking technology ready to tackle contemporary challenges. This is part of the ``rebranding Prolog'' effort that we believe is essential.




\section{Extending Prolog's user base}\label{sec:popularity-factors}

\begin{quotation}
\noindent There are people who don't like popularity. It's much better to be exclusive and remote.
\flushright\textit{Robert Indiana}
\end{quotation}

\subsection{The imperative for growth}

Prolog \textit{used} to be somewhat popular. In the eighties, during the second AI summer, the Japanese Fifth Generation Computer Systems (FGCS) project placed Prolog at the center of the computing world. Communities formed, excellent textbooks were written, and the language became more efficient and expressive.

Today, the landscape is different. The authors of \textit{Fifty Years of Prolog and Beyond} (FYPB) admit that Prolog has a ``comparatively small user base.'' This is an understatement. Data presented by Statista estimates that in 2023 there were around 28 million programmers worldwide. By contrast, the SWI-Prolog forum -- likely the largest gathering of Prolog developers -- has only around 1,300 members. The Erlang forum has roughly 2,300. Compare this to Python, which boasts over 1,600 \textit{local} user groups across 191 cities. As far as we are aware, there are no local Prolog user groups in any city.

The authors of FYPB identify this comparatively small user base as a significant threat. The Prolog community faces a crisis not of capability, but of invisibility. It is not just that Prolog is underrepresented in educational curricula or industry job postings. More critically, it lacks the mainstream discourse presence and cultural visibility that define a healthy language ecosystem. Few conferences center on Prolog, and few young programmers encounter it outside of narrow academic contexts.

Why does this matter? Prolog is an excellent language, but being good does not by itself make a language popular. In fact, the causality often works the other way around: a popular language is more likely to improve and become better. Popularity is not merely an intrinsic merit or a vanity metric; it is the engine of ecosystem health. Widespread use brings crucial benefits: more libraries, better tooling, more critical feedback, and greater pressure for innovation. Paul Graham articulated this dynamic in \textit{Hackers \& Painters} (2004):

\begin{quote} So whether or not a language has to be good to be popular, I think a language has to be popular to be good. And it has to stay popular to stay good. The state of the art in programming languages doesn't stand still. And yet the Lisps we have today are still pretty much what they had at MIT in the mid-1980s, because that's the last time Lisp had a sufficiently large and demanding user base. \end{quote}

\noindent While the Lisp ecosystem has since seen a revival through Clojure, Graham's central point holds true: a large user base fuels the ecosystem that, in turn, increases the utility and appeal of the language itself.

%
%
%\noindent %Created in 1972, Prolog is a rather old programming language. Any rumors of its imminent death can easily be debunked, however, as it is currently experiencing something of a revival. New systems are being implemented, features such as tabling and constraint logic programming libraries are being added to new as well as already established systems, new books are being written, the community is small but thriving, and the language is still taught at some universities. (As already mentioned, the paper \textit{50 Years of Prolog and Beyond} \cite{??} gives an excellent view of the evolution of Prolog during its 50 years of existence.)
%Prolog \textit{used} to be somewhat popular, but that was in the eighties, when the second AI summer was in full swing, when the Japanese ran the Fifth Generation Computer Systems (FGCS) project, when most of the (still very) good textbooks were being written, and the communities formed that greatly improved the Prolog language by making it more efficient and more expressive.
%
%The authors of \textit{Fifty Years of Prolog and Beyond} admit that Prolog is ``not a mainstream language'' and that it has a ``comparatively small user base.'' The phrases quoted are clearly understatements. Data presented by Statista estimates that in 2023 there are around 28 million programmers worldwide.\footnote{\url{https://www.statista.com/statistics/627312/worldwide-developer-population/}} Let us be optimistic here, and assume that a couple of thousand developers are using Prolog. At the time of writing, the SWI-Prolog forum\footnote{\url{https://swi-prolog.discourse.group}} has 1053 members, and this is most likely the largest Prolog forum. The Erlang forum has 1435 members. Compare this with Python for example, which has about 1,637 \textit{local} Python user groups worldwide in almost 191 cities, 37 countries and over 860,333 members.\footnote{\url{https://wiki.python.org/moin/LocalUserGroups}} As far as we are aware, there are no local Prolog user groups, in any city.
%
%The authors of \textit{Fifty Years of Prolog and Beyond} (FYPB) sees the comparatively small user base as a threat to Prolog. The Prolog community faces a crisis not of capability, but of invisibility. It is not just that Prolog is underrepresented in educational curricula or tech industry job postings. More critically, it lacks the local user groups, the mainstream discourse presence, and the cultural visibility that define a healthy, growing language ecosystem. Few conferences center on Prolog, and few young programmers encounter it outside of narrow academic contexts. 
%
%Prolog is an excellent language, but being good does not by itself make a language popular. In fact, the causality tends to work the other way around; a popular language is more likely to be improved and become better. In the words of Paul Graham in \cite{graham2004hackers}:
%
%\begin{quote}
%So whether or not a language has to be good to be popular, I think a language has to be popular to be good. And it has to stay popular to stay good. The state of the art in programming languages doesn't stand still. And yet the Lisps we have today are still pretty much what they had at MIT in the mid-1980s, because that's the last time Lisp had a sufficiently large and demanding user base.
%\end{quote}
%
%\noindent Now, 15 years later, it is debatable whether Graham was correct about Lisp. After all, the Clojure programming language, which appeared some fifteen years ago, improved Lisp  (most would say). Also, a number of logic programming researchers and Prolog implementers have managed to greatly improve Prolog despite its lack of popularity, presumably because its use in research and teaching.
%
%However, we believe that Graham's main point still holds. Popularity is not an intrinsic merit, but widespread use brings crucial benefits: more libraries, better tooling, more critical feedback, and greater pressure for innovation. A larger user base fuels an ecosystem that, in turn, increases the utility and appeal of the language itself. In this view, the Prolog community's small size is not just a demographic detail -- it is a technical liability.
%
%So in order to make Prolog an even better language, it makes sense to try to make it more popular. so how do we do that?\footnote{Another argument for why it would be great if Prolog became more popular is that a greater popularity leads to more Prolog code being written, something that allow systems of Generative AI to train on more code, improving the help one can get from them.} Unfortunately, the authors of FYPB do not have much to say about it.
%
%---

%\subsection{What makes Prolog worthy of revival?}
%
%Prolog's decline in visibility is not the result of inherent shortcomings. Quite the opposite: it offers a number of unique strengths that continue to make it a compelling choice in specific domains. It embodies a paradigm -- logic programming -- that aligns closely with symbolic AI, formal reasoning, and declarative knowledge representation.
%
%With its actor-style concurrency, message-passing architecture, and support for network-transparent computation, Web Prolog brings modern distributed programming idioms into the logic programming fold. It bridges the gap between symbolic reasoning and the real-time, agent-oriented logic of the Web.
%
%These are not minor upgrades. They amount to a reimagining of what Prolog can be -- not just a teaching language or a theorem prover, but a foundation for a new kind of Web programming.

\subsection{Key factors for popularity}

Graham has more to say about popularity and what it takes for a programming language to become popular:\footnote{\url{http://www.paulgraham.com/popular.html}}

\begin{quote}
Programming languages don't exist in isolation. To hack is a transitive verb -- hackers are usually hacking something -- and in practice languages are judged relative to whatever they're used to hack. So if you want to design a popular language, you either have to supply more than a language, or you have to design your language to replace the scripting language of some existing system. 
\end{quote}

\noindent Applying this logic, supplying an entire web-based ecosystem, rather than just a language in isolation, presents a compelling path forward and offers a viable alternative that could replace JavaScript for some applications within this specific domain.

Antonio Cangiano, a software developer and technical evangelist at IBM, provides a more comprehensive list of key factors for popularity to consider:\footnote{\url{http://programmingzen.com/so-you-want-your-programming-language-to-be-popular/}}

\begin{enumerate}
 
\item Being a default language for a popular ecosystem;

\item Being backed and promoted by a tech giant;

\item Having a large standard library and/or targeting a popular VM (e.g., JVM), unless it is a system language;

\item Fully embodying a new paradigm shift;

\item Being very useful in a particular domain, like Ruby did through Rails back in 2005;

\item Having great documentation, guidance for newcomers, tools, editor integration, and in general removing any instrumentation obstacle that makes getting started with your language difficult;

\item Fostering a welcoming community that encourages sharing, marketing, and evangelism;

\item Having a somewhat familiar syntax that is easy to read, write, and teach;

\item Being active and developed for several years;

\item Providing technical innovations that lead to productivity and more maintainable code;

\item Luck.

\end{enumerate}

\noindent As Cangiano points out, the first two key factors are out of the question for most languages, and even then, few languages will meet all of the remaining requirements. Even embracing just a few of these should, however, be enough to grant a language a chance at becoming mainstream. So how well does Prolog satisfy these requirements, and how can Web Prolog help?

As for being the default language for a popular ecosystem (e.g. JavaScript for the Web), we are trying to \textit{build} an ecosystem, namely the Prolog Trinity, for which Web Prolog is meant to be the default language. Whether the Prolog Trinity ecosystem will become popular remains of course to be seen. But we do not have to rely on the entire Prolog Trinity becoming a success. Even if the Prolog Web only finds few uses, Web Prolog as a special-purpose web programming profile of Prolog may still have an interesting role to play.

Being backed by a tech giant would certainly be nice, but being backed by the W3C might also help. If standardization is successful, industry backing and corporate sponsors that drives marketing or integration work might well show up.

Web Prolog does not target a specific VM. Any VM or any programming language that can possibly support an implementation of a profile of Web Prolog is targeted. After all, this is the whole idea behind a language aiming at real-time interoperability between different programming systems that may have different VMs or no VM at all.

As for fully embodying a new paradigm shift, we could possibly claim that the combination of web logic programming and web agent programming proposed in this book might be counted as a new paradigm, based on the two older paradigms of logic programming and actor programming. However, it may be more reasonable to think of Web Prolog as a multi-paradigm programming language.

Regarding the usefulness in a particular domain, creating a domain-specific profile of Prolog is our main focus -- we aim to make Web Prolog very useful in the web domain. To aim for the kind of dominance Ruby on Rails had in this domain would be shooting for the stars. In addition, optionally in combination with SCXML, Web Prolog also tries to become a viable language for implementing sophisticated web-based agent programming platforms and intelligent conversational systems.

Documentation of the Web Prolog built-in predicates can be pooled from the documentation provided by existing Prolog systems. Guidance for newcomers in the form of textbooks, many of which are free, exists. As for removing other obstacles that make getting started with Web Prolog difficult, web-based IDEs such as SWISH must again be mentioned.

Currently, the most active Prolog community is arguably the one formed around SWI-Prolog. It is indeed a welcoming and encouraging community.

Whether or not having a somewhat familiar syntax that is easy to read, write and teach applies to Web Prolog depends on where you are coming from. For developers raised on Python or JavaScript the syntax of Prolog may seem alien at first, but for an Erlang programmer, Web Prolog will present few problems. %More on this will be said in Section~\ref{sec:learning-web-prolog}.

Regarding the factor of ``having been active and developed for several years,'' we can of course not claim this for Web Prolog. However, the first Prolog system was developed already in the early seventies, and Erlang was conceived of and implemented during the late nineties, so from this perspective Web Prolog certainly has roots stretching deep into the history of programming languages, and through shared DNA with its parent languages a certain maturity of the underlying ideas can be claimed for it.

As for providing technical innovations leading to productivity and more maintainable code, this is what Prolog and technologies such as constraint logic programming is for.

And yes, we would certainly need luck as well, but as the Spanish proverb says, luck comes to those who look after it. 

To these key factors, we are going to add another one, namely the existence of a standard. For a web technology such as Web Prolog, where interoperability between different implementations is at stake, standardization is particularly important. %We shall spend the whole of Chapter~\ref{sec:standardisation} discussing the prospect of a new standardization effort.\\

Standardized languages tend to have larger communities and ecosystems surrounding them. This means there are more resources available, such as documentation, tutorials, libraries, and frameworks, which can accelerate development and provide solutions to common challenges. A robust community fosters collaboration, innovation, and knowledge sharing.

%\subsection{Ways to make Prolog popular again}\label{}
%
%\begin{quotation}
%\noindent There are people who don't like popularity. It's much better to be exclusive and remote.
%\flushright\textit{Robert Indiana}
%\vspace{0.5cm}
%\end{quotation}
%
%\noindent The wordings of the heading of this chapter seem to make the (in our opinion) false presupposition that Prolog is not already in many ways a great programming language. But if greatness of a programming language is measured in number of practitioners, then the presupposition is, unfortunately, true. In this chapter we will take a stab at the problem of popularising Web Prolog. We have quite some way to go. Estimates have it that there are 18.5 million software developers in the world (including ~7.5 million hobbyists).\footnote{\url{https://www.infoq.com/news/2014/01/IDC-software-developers}} Let us be optimistic here, and assume that a couple of thousand developers are using Prolog. This is nowhere near greatness.
%
%---
%
%Also note that while overall the number of software developers have grown dramatically since the mid-nineties, the number of Prolog developers have fallen. However, the fact that the language refuses to die should tell us something and give us hope for a future revival.
%
%When we guess that there may be perhaps a couple of thousand Prolog programmers in the world, we do not count people using Prolog as part of their education in university computer science courses. It is of course a good thing that teachers still see Prolog as interesting enough to teach it for a couple of weeks in order to introduce a new and different paradigm, but students studying Prolog in a four-languages-in-eight-weeks course should not be counted as Prolog developers. 
%
%Most students are not capable of learning the language is such a short time. Is it too difficult to learn? While it may be true not all people are wired for Prolog, we do think some are. Likely, they can be found among the very best of JavaScript programmers. Therefore, it might be a good idea to actively seek out and educate these people, showing them what can be done with a combination of JavaScript and Web Prolog. [FIXME: Move this paragraph somewhere else.]
%
%===
%
%\noindent Paul Graham has written about popularity too:\footnote{\url{http://www.paulgraham.com/popular.html}}
%
%\begin{quote}
% Programming languages don't exist in isolation. To hack is a transitive verb -- hackers are usually hacking something -- and in practice languages are judged relative to whatever they're used to hack. So if you want to design a popular language, you either have to supply more than a language, or you have to design your language to replace the scripting language of some existing system. -- \textit{Paul Graham}
%\end{quote}

%\subsection{Marketing}\label{sec:marketing}
%
%We asked OpenAI's ChatGPT to come up with some marketing text, and this is what it produced:
%
%\textbf{Web Prolog} is Prolog re-imagined for an Internet-scale, actor-oriented world: a lean ISO-compatible core fused with Erlang-style message passing so that every term can travel, every goal may be remote, and every process owns a private mailbox; logic variables, unification and backtracking yield declarative power, while links, monitors and supervisors deliver fault-tolerant concurrency; and first-class HTTP and WebSocket predicates give native web reach, enabling the same program to run in the browser, on the server or at the edge and to glue together classic Prolog engines, reasoners and knowledge stores -- think Erlang's robustness, Prolog's reasoning and the Web's connectivity rolled into a single, potentially standardizable \textit{lingua franca} for logic-based systems.
%
%\textbf{Prolog agents} are millisecond-spawned, logic-driven actors that perceive, reason and act across the Web: each agent is simply a Web Prolog process with a mailbox, private knowledge base and conversational protocol, so anything from a chatbot to a crawler to an intelligent conversational agent speaks pure Prolog terms; agents scale from toy helpers to long-running services, can supervise or be supervised and compose into multi-agent organisations -- nodes, toplevel actors, statechart actors, or your own specialities -- and because their ``brains'' are Prolog they provide built-in pattern matching, rule-based inference, natural-language grammars and explainable traces, thereby turning declarative knowledge into live, autonomous web services that you can spawn, link and converse with as easily as you assert a fact.
%
%\textbf{The Prolog Web} is a distributed fabric in which code and knowledge flow freely among logic-aware nodes: each node exposes HTTP and WebSocket endpoints, accepts Web Prolog code and hosts agents that answer queries with proofs rather than strings; send/receive works almost transparently whether the target is next door or two continents away, yet security, naming and fault domains remain explicit and enforceable; knowledge can travel to computation or computation to knowledge, enabling data-sovereign analytics, edge reasoning and collaborative AI; and seamless interoperability with JSON, RDF and other web formats adds an executable logic layer atop the Semantic Web -- imagine today's Web, but every URI can explain itself because a Prolog agent behind it is ready to justify its answers.



\subsection{Rebranding Prolog}\label{sec:rebranding-prolog}

\begin{quotation}
\noindent \textit{Rebranding} is a marketing strategy in which a new name, term, symbol, design, or combination thereof is created for an established brand with the intention of developing a new, differentiated identity in the minds of consumers, investors, competitors, and other stakeholders.
\vspace{-0.3cm}
\flushright\textit{Wikipedia}
\vspace{0.5cm}
\end{quotation}

\noindent %\textit{Strategic ambiguity} is a marketing trick designed to make customers think they are getting one thing when in reality they are getting something else. In our case we offer people shopping for a programming language that suits their needs the \textit{appearance} of a new language. It is a trick, but is it \textit{dirty}?
Here is an interesting quote from a Hacker News user that calls himself ``gruseom:''\footnote{\url{https://news.ycombinator.com/item?id=3900047}},\footnote{His real name is Daniel Gackle and he is actually the head of Hacker News.}

\begin{quote}
With programming languages, I've come to the conclusion that ``If X is so awesome, why is nobody using it?'' is the wrong question to ask. Or rather, it's a fair question, but the answer doesn't depend on X. It is overwhelmingly a matter of network effects and social proof. What matters is what everyone else is using.

A programming language has a window of opportunity to gain mindshare while it's new. Once that window passes, the odds are hugely against it. It's true there are a few technical factors, for example that older implementations lack the infrastructure to work with newer technologies, but such things can be improved. The dominant factor, the real barrier, is social psychology. This is obscured by our tendency to retrofit plausible-sounding reasons to our choices (``X is slow'' or ``it doesn't have libraries'' or whatever).

This does suggest a strategy for reviving an old language: rebranding. First, you need a new implementation. Trying to convince people to use the old one is like trying to get them to watch an old movie. George Lucas didn't promote Kurosawa, he made Star Wars. Second, there has to be a hook, something to convince us that this really is a new language, modern and with the times. So, light sabers. That's critical for opening a fresh window of opportunity instead of getting pegged to the old, long-closed one. It needn't be the most important thing, but it can't be fake either; it must be real enough or no one will take the new brand seriously. Without a convincing reason for why we weren't using X before, the new window will never open. Once it opens, then and only then do the beauty and power of X get their chance to bond with the user. In the Star Wars analogy, that would be the characters and the story -- the deeper things that came from Kurosawa and ultimately from ancient archetypes. They'd never have gotten the chance to kick in if Star Wars hadn't got the audience in the theater. Once they did, though, then you had Star Wars fans for life.

With Clojure, the hook was Java libraries. Actually, there were three hooks. The others were concurrency and sequences. But Java libraries is the one that worked. And now there's a modern Lisp again!
\end{quote}

\noindent With Web Prolog, the hook is the Prolog Web with its focus on bringing more logic to the Web and the promise of real-time interoperability between different Prolog (and non-Prolog) systems. Not every new or rebranded programming language comes with a novel extension of the Web, so from a marketing perspective it would make sense to emphasize the Prolog Web as the ``killer application''  (or ``killer middleware''?) for Web Prolog. Actually, there are three hooks here too. The other hooks are actors and a mature concurrency programming model with a proven track record thanks to Erlang and Elixir. 

Since the author quoted above regards Clojure as a rebranded Lisp, and since it seems reasonable to regard Elixir as a rebranded Erlang, a comparison between the three ``rebranding projects'' may be interesting. Figure~\ref{fig:clojure-comparison} is meant to illustrate such a comparison.

\begin{figure}[h]
    \centering
    \includegraphics[width=9cm]{clojure-vs-erlang-vs-web-prolog.png}
    \caption{Lisp rebranded into Clojure and Erlang rebranded into Elixir, compared with Prolog rebranded into Web Prolog}
    \label{fig:clojure-comparison}
\end{figure}

\noindent Apart from the fact that we \textit{know} that Lisp was successfully rebranded into Clojure, and Erlang into Elixir, but are less sure of how it would work out in the case of Web Prolog as a rebranded version of Prolog, there are other notable differences. In contrast to general-purpose languages such as Clojure end Elixir, Web Prolog is a \textit{profile}, not a fork. For this reason, also in contrast to Clojure end Elixir, Web Prolog places great weight on backwards compatibility, i.e. on being able to run the vast majority of examples appearing in old but still very good Prolog textbooks. In other words, for Clojure and Elixir, the language has changed, but not its purpose. For Web Prolog, the language will stay (almost) the same, but its purpose has become a different one, namely web programming.  

Moreover, Web Prolog does not run on top of a VM, but on top of the Web, propped up by nodes that serve as runtime systems disguised as web servers. This makes a lot of sense for a web programming language. Also, while the motive behind the rebranding of Lisp may well have been to create (what ``gruseom'' refers to as) a \textit{modern} language, this is not the motivation behind Web Prolog. We are happy with ``good ol' Prolog'' adapted to the Web, and with a standard that would allow us to regard Web Prolog as a web technology in its own right.


The history of Elixir offers a specific lesson for us. Elixir did not become popular  solely on its language merits; it exploded in popularity because of the Phoenix web framework. Phoenix provided a ``killer app'' environment -- a high-performance,  real-time web framework that made the power of the Erlang VM accessible and  productive for web developers.  The Prolog Trinity aims to create a similar effect. By providing a complete  ecosystem -- the Prolog agents and the Prolog Web architecture -- we offer more than  just a language; we offer a platform. Just as Phoenix made Elixir the go-to choice  for scalable web apps, the Prolog Trinity can make Web Prolog the go-to choice for  intelligent, distributed agents.  

%The discussions throughout this chapter implicitly and explicitly point towards a critical objective: the significant expansion of Prolog's user base and its sphere of influence. This is not a goal pursued for reasons of vanity or mere popularity metrics; it is an existential imperative for the long-term health, vibrancy, and continued evolution of Prolog and the broader Logic Programming ecosystem.

The ultimate purpose of this rebranding effort, and of the broader work of constructing an ecosystem around Web Prolog, is to set a virtuous cycle of adoption in motion. As the community grows, the ecosystem becomes richer: more users generate more libraries, more tools, and more opportunities for integration. A larger user base also produces a broader and more accessible corpus of learning materials -- tutorials, courses, examples, and best-practice guides -- which directly addresses one of Prolog's long-standing challenges, namely its reputation for a steep learning curve. With broader engagement comes attention from industry, and with that attention follows investment, commercial experimentation, and a willingness to fund further research and development.

The notion of ``rebranding Prolog'' is central to this strategy. It is not about hiding Prolog's identity, but about reshaping its narrative. We are moving from a story about the past -- expert systems and 1980s AI -- to a story about the future: a web of intelligent, interoperable agents. Extending the user base is the only way to ensure that Prolog not only survives but thrives, contributing its unique logic programming paradigm to the next fifty years of computing.




%\begin{itemize}
%   \item \textbf{The Virtuous Cycle of Adoption:} Popularity in the world of programming languages and technologies creates a powerful virtuous cycle:
%   \begin{itemize}
%       \item \textbf{Richer Ecosystems:} A larger user base naturally leads to the development of more libraries, frameworks, tools, and integrations, making the language more productive and versatile for everyone.
%       \item \textbf{Abundant Learning Resources:} More users stimulate the creation of diverse and high-quality learning materials -- books, tutorials, online courses, sample projects, and community forums -- lowering the barrier to entry for newcomers.
%       \item \textbf{Stronger Community Support:} A thriving community provides more avenues for collaboration, mentorship, problem-solving, and the sharing of best practices.
%       \item \textbf{Increased Industry Demand and Job Opportunities:} Wider adoption in industry translates into more projects utilizing the technology, creating demand for skilled developers and, consequently, more job opportunities.
%       \item \textbf{Greater Investment and Innovation:} Success attracts investment, both from commercial entities and academic institutions, fueling further research, development, and refinement of the language and its supporting infrastructure.
%   \end{itemize}
%   \item \textbf{Overcoming Perceived Niche Status:} Despite its undeniable power and elegance, Prolog has often been perceived as a ``niche'' language, primarily confined to academic research or specialized AI applications of a bygone era. This perception, whether entirely fair or not, can be a self-fulfilling prophecy, discouraging potential new users and limiting its consideration for a broader range of projects. The Prolog Trinity is a deliberate strategy to break free from these confines, demonstrating Prolog's applicability and superiority in thoroughly modern domains like web services, concurrent programming, and intelligent agent systems.
%   \item \textbf{Key Factors Influencing Adoption -- And How the Trinity Addresses Them:} We can analyze language adoption through several lenses, and the Prolog Trinity is designed to make positive inroads on many fronts:
%   \begin{itemize}
%       \item \textbf{Default Language for a Compelling Ecosystem:} Just as Swift is to iOS or C\# to .NET, Web Prolog aims to be the natural, indispensable language for the Prolog Trinity ecosystem -- a rich environment of interoperating logical agents and services.
%       \item \textbf{Solving Real-World Problems Better:} Ultimately, developers adopt tools that help them solve their problems more effectively. The Trinity targets problem domains (complex rules, distributed intelligence, semantic integration) where logic programming offers distinct advantages.
%       \item \textbf{Lowering the Barrier to Entry:} Through better tools, documentation, ``Prolog Playgrounds,'' and a focus on web-friendliness, we aim to make Prolog more accessible.
%       \item \textbf{Fostering a Vibrant and Welcoming Community:} The Prolog Web concept is as much about social infrastructure as technical infrastructure.
%       \item \textbf{Demonstrating Continuous Innovation:} The Prolog Trinity itself is a significant innovation, signaling that Prolog is not a static language but one that continues to evolve and adapt.
%   \end{itemize}
%   \item \textbf{The ``Rebranding'' Imperative Revisited:} The notion of ``rebranding Prolog'' is central to this growth strategy. It involves actively shaping a new narrative for Prolog -- one that emphasizes its dynamism, its web-savviness, its unique strengths in building intelligent and concurrent systems, and its readiness to meet the challenges of the 21st century. This is about highlighting relevance and future potential, not just past glories.
%\end{itemize}
%Extending the user base is the pathway to ensuring that Prolog not only survives but thrives, continuing to contribute its unique logical paradigm to the art and science of computation for generations to come.

  

  
%\section{A plan, and a project}


\section{Risk analysis}\label{sec:path}

\begin{quotation}
I've heard that the Web is dead. But all the applications that have killed it are accessing services using HTTP over port 80.  
\vspace{-0.3cm}
\flushright\textit{Nat Torkington}\footnote{As quoted by Mike Loukides in ??, attributing the remark to Nat Torkington.}
\end{quotation}

\noindent In a thread on SWI-Prolog's Discourse forum addressing the ``marketing'' of Prolog, Paulo Moura voiced some scepticism regarding our approach:\footnote{\url{https://swi-prolog.discourse.group/t/encouraging-industry-about-prolog/1106/58}}

\begin{quotation}
As important as the Web is, and therefore the significance of Prolog being a good programming language for Web applications, the rebranding you suggest is overreaching and doesn't make sense to me. Prolog, and logic programming in general, is much more than a specific application field. It's already bad that, in some circles, Prolog is being downplayed as a DSL tool. The last thing we need is for Prolog to be downplayed as a Web thing. Prolog is a general programming language and that's the vision a lot of people in the community is working hard to establish. Trying to sell Prolog was a Web silver bullet would only backfire in the same way Prolog association with the hype surrounding AI backfired in the past. While the Web can be a stage to increase logic programming marketshare and mindshare, is not and should not be, in my view, the end game.
\end{quotation}

\noindent %It may have been a mistake to refer to Web Prolog as a ``special-purpose language'' and advocating ``rebranding.''  
Upon reflection, describing Web Prolog as a ``special-purpose language'' was too narrow. Prolog should be recognised as a general-purpose language, while Web Prolog is better characterised as a specialised \textit{profile} -- a carefully delimited subset of the larger Prolog language. The right way to state the intent is therefore not that Prolog should be redefined as a web language, but that it should gain a clear web-facing profile and an interoperable surface. Prolog remains a general-purpose language; Web Prolog is a specialised profile that makes certain choices explicit at the boundary, so that nodes and clients can interoperate across the open network. The Web, on this view, is not an end game or a silver bullet, but a particularly effective stage on which Prolog systems can be made discoverable, composable, and runnable by others.

The deeper point is that ``web-facing'' need not mean ``web-only.'' Many general-purpose languages became widely used by first becoming indispensable in a particular arena: SQL in databases, JavaScript in browsers, R in statistics. A credible entry point is not a loss of generality; it is often how generality is earned.

The Web is a particularly plausible entry point for Prolog because many web problems have an implicit logical shape: query and transformation over structured data, policy and access control, constraint-based selection, and reasoning over graphs. Presenting Prolog as a rules-and-constraints component inside a polyglot web architecture is therefore not a marketing trick; it is an alignment between strengths and an enormous application surface.

The main failure mode to avoid is not ``being on the Web,'' but making vague claims. The Trinity's stance is deliberately modest: specify clear interfaces, provide runnable nodes, and let adoption be driven by concrete utility rather than by rhetoric.

%Moreover, languages rarely win mind-share by advertising themselves as general, all-purpose tools. SQL is a query language, JavaScript a browser language, R a statistics language -- and yet all three have long since escaped those boundaries.  What conferred legitimacy was a first conspicuous success in a domain that matched the language's intrinsic strengths.  Giving Prolog a similarly sharp public image in the web arena is therefore not a down-play but a proven adoption tactic.
%
%Modern web stacks wrestle with data representation, access-control policies, constraint-based scheduling, and knowledge-graph querying.  All of these boil down to search, unification, and back-tracking -- exactly Prolog's core competencies.  Positioning Prolog as the declarative engine inside a polyglot web architecture is technologically coherent, not opportunistic.
%
%The AI winter was triggered less by publicity itself than by unbounded promises.  A disciplined claim -- ``Prolog is your rules engine for complex web data flows'' -- is concrete, verifiable, and modest compared to the grand, ill-defined ``AI language'' label of the past. Specificity is an antidote to hype, not its synonym.
%
%Treating Prolog as ``a Web thing'' is not a surrender of generality but a pragmatic market-entry strategy: it aligns the language with an enormous, rapidly evolving problem space that naturally rewards declarative reasoning, leverages existing mature tooling, and offers maximal visibility to new developers.  Far from backfiring, such focus is the most realistic path to the broader recognition that many in the community have long desired.

Wielemaker agreed with Moura, and followed up with the following statement:

\begin{quotation}
Prolog is quite suitable for the web, but in my experience Prolog is used for a very diverse set of tasks, many of which need Prolog more than the web. I'd love to see Web Prolog fly, in which case it is not unlikely to become more or less detached from the Prolog systems we see targeting other applications.
\end{quotation}

\noindent While we do agree that Prolog is used for many different task not related to the Web, we believe that there are ways to ensure that Web Prolog does not become detached from systems targeting other applications. The key here is to regard Web Prolog, not as a special-purpose language, but as a \textit{profile} of the general-purpose Prolog language, and as a profile of a \textit{future} Prolog. We refer to this future language as ISO Prolog 2.0.

We have a proposal for a path for Prolog that takes us from ISO Prolog to Web Prolog and then from Web Prolog to ISO Prolog 2.0. Such a path can be discerned against the backdrop of the profile diagram in Figure~\ref{fig:profile-hierarchy-full}, which extends the diagram introduced in Chapter~\ref{sec:prolog-agents} with yet another profile, for ISO Prolog 2.0. 

\begin{figure}[h]
    \centering
	\includegraphics[width=\textwidth]{profile-hierarchy-full}
    \vspace{0.3cm}\caption{Web Prolog as a subset of a hypothetical future ISO Prolog 2.0.}
    \label{fig:profile-hierarchy-full}
\end{figure}

ISO Prolog 2.0 might include (at least) what Web Prolog did not include from ISO/IEC 13211-1:1995, such as predicates for file I/O. It might also include tabling, constraint logic programming or whatever technologies research into logic programming has come up with when the world is ready for it. Perhaps it might even be the Unified Prolog discussed earlier in this chapter.

Our proposal means, and it follows from the Venn diagram, that an implementation of an ISO Prolog 2.0 system would only be conformant if it implements the Erlang-style predicates for concurrency and distribution as well as the web APIs on which distribution depends. This might involve other forms of APIs as well (such as raw TCP/IP) and an extended set of predicate options, that allow Erlang-style distribution in a closed environment to work.

This suggests that a plan for the design and implementation of Web Prolog can easily be made to ``fit inside'' a plan for the design and implementation of ISO Prolog 2.0, one project inside another.

It so happens that concurrency is included in the recipe for a Unified Prolog. Could the community be persuaded that Erlang-style concurrency is the way to go, then standardizing the relevant built-in predicates can be commenced already now, which basically means that the design and development of Web Prolog and of ISO Prolog 2.0 can be done in parallel.


One hope that we may have is that Web Prolog can inspire more and intensified work on the ISO Prolog standard as well as on the Prolog Commons. Perhaps the standardization of constraint logic programming is the next task for the ISO working group to tackle? Or perhaps finalizing DCG is more important.
